\section{Related Work}

This section reviews the research and engineering practices needed to turn LLM-agent prototypes into traceable enterprise systems. The organizing question is how an answer can be grounded, assembled, constrained, and validated before it is presented to the user.

\paragraph{RAG and traceability.}
Retrieval-augmented generation (RAG) combines language models with external documents to improve factual grounding and updateability \citep{lewis2020rag,gao2023rag}. Later work shows that retrieval must also be selective and checkable: self-reflective retrieval methods ask whether retrieval is needed and whether generation is supported \citep{asai2024selfrag}, while factuality metrics decompose long-form answers into atomic claims \citep{min2023factscore}. Retrieval selectivity and factual checkability are directly relevant to enterprise domains where disclosures, investor-relations (IR) materials, news, and regulatory filings change over time. Retrieval provides factual context, while traceability additionally requires source state, entity scope, and claim-level support. Retrieved passages may be stale, irrelevant, or mixed across entities, and fluent generation may still hallucinate \citep{ji2023hallucination}. Our approach treats RAG as one component inside a structured source-to-claim harness.

\paragraph{LLM agents and orchestration.}
Classical agent research emphasizes autonomy, interaction, and explicit system architecture \citep{wooldridge1995intelligent,jennings1998roadmap}. LLM agents extend the classical agent tradition through tool use, planning, search, and intermediate execution traces \citep{yao2023react,schick2023toolformer,wang2024surveyagents,acharya2025agentic}. Tool-use benchmarks further show that API selection and execution are engineering concerns as well as prompting concerns \citep{qin2024toolllm}. Recent applied work also reports multi-platform agent deployment and programmable agent runtime architectures \citep{ahn2025autonomous,walters2025eliza}. The present paper shifts the emphasis from multi-platform deployment to enterprise hardening: each answer must record which sources, tools, claims, and fallback paths shaped the final output.

Recent orchestration frameworks make agent engineering more concrete. AutoGen supports multi-agent conversation and coordination \citep{wu2024autogen}; LangChain Agents provide prebuilt tool-calling loops and LangGraph provides graph-based runtimes for stateful agents \citep{langchain2026agents,langchain2026langgraph}; and CrewAI organizes role-based agents, crews, and flows \citep{crewai2026software}. The proposed control layer is complementary to these frameworks: orchestration frameworks address how agents are composed and executed, while the pattern studied here addresses whether each produced answer is grounded in source manifests, claim records, output contracts, and validation traces. In deployment, it can surround a LangGraph, CrewAI, AutoGen, or LangChain-agent workflow.

\paragraph{Prompt engineering and prototyping.}
Prompt engineering became important because large language models can perform new tasks from natural-language instructions and in-context examples \citep{brown2020language,liu2023pretrain}. Prompt engineering also entered broader academic and practitioner discourse as users learned to shape model outputs through prompt design \citep{giray2023prompt}. Techniques such as chain-of-thought prompting further showed that changing the prompt can substantially change reasoning behavior \citep{wei2022chain}. Vibe coding extends prompt design from isolated task instructions to an iterative development workflow in which software behavior is co-created through natural-language dialogue with AI tools \citep{karpathy2025vibecoding,meske2025vibecoding,malamas2025vibecoding}. Prompting and vibe-coding techniques are valuable for early exploration and are one reason prototype systems can be built quickly. The transition addressed by the present paper appears when prompts stop being only task instructions and become the place where source policy, entity routing, answer structure, data freshness rules, and user-interface behavior are implicitly stored. At that point, the system becomes prompt-dominant: effective for exploration, while audit, replay, and transfer require stable responsibilities to be moved into explicit artifacts.

\paragraph{Programmable LLM pipelines.}
Several frameworks argue that LLM applications should be programmed, tested, and optimized with explicit software abstractions. Demonstrate-Search-Predict and DSPy represent language-model pipelines as composable declarative programs \citep{khattab2022demonstrate,khattab2024dspy}; the DSPy project repository is also cited as a software artifact \citep{stanfordnlp2026dspysoftware}. LMQL and Guidance treat model interaction as controlled execution \citep{beurerkellner2023lmql,guidance2026software}; and Instructor uses typed schemas to constrain outputs \citep{instructor2026software}. Runtime guardrail systems such as NeMo Guardrails similarly show how user-defined rails can control conversational behavior independently of the underlying model \citep{rebedea2023nemo}. Autoresearch illustrates how improvement can be organized as repeatable experiments with explicit logs \citep{karpathy2026autoresearch}. Our contribution is to apply the programming-oriented philosophy at the product-architecture level: routing, source eligibility, claim promotion, answer planning, follow-up filtering, and internal-output leakage prevention are implemented as code and data structures.

\paragraph{LLM Wiki and knowledge management.}
Karpathy's LLM Wiki pattern proposes a maintained markdown knowledge layer between raw sources and query-time answers \citep{karpathy2026llmwiki}. In that pattern, the LLM helps maintain a persistent wiki of entity pages, topic summaries, links, contradictions, and logs, reducing repeated synthesis from the same raw documents. The pattern is useful for research and business settings where knowledge accumulates over time. The proposed system adopts the maintained-knowledge idea while assigning authority to source manifests and source-backed claims: raw source manifests and source-backed claims remain the source of truth, and the LLM Wiki is compiled as a concise context layer for humans and models.

\paragraph{System validation and evaluation.}
Software-engineering research on AI-based systems emphasizes that production AI depends on data management, testing, lifecycle controls, monitoring, and assurance artifacts in addition to model performance \citep{martinezfernandez2022seai,ashmore2021assuring}. Machine learning operations (MLOps) research similarly treats operationalization as a lifecycle concern that spans the full system lifecycle \citep{kreuzberger2023mlops}. RAG evaluation frameworks measure faithfulness, answer relevance, context precision, and context recall \citep{es2023ragas,saadfalcon2024ares}, while broader benchmark work supports multi-scenario evaluation \citep{liang2023helm}. Practitioner tools also support prompt, model, and application regression checks \citep{promptfoo2026software}. The present paper uses system validation suited to a pre-commercial research prototype before customer logs are available. The focus is whether fixed scenarios preserve source-backed claims, entity routing, trace contracts, answer structure, interface status, and latency.
